\section*{Potenzreihen Aufgaben mit L�sungen}\label{MA2-potenzreihen-loesungen}

\ifintegriereaufgaben Erkl�rungen und Beispiele, die zum Verst�ndnis der L�sungswege hilfreich sind, finden Sie im Skript\wennreferenzvorhanden{MA2-potenzreihen}{ (Abschnitt \abschnittsverweis{MA2-potenzreihen})}. \fi

\begin{exercisebox}[Konvergenz] %Papula VI, Abschnitt 2, Aufgabe 1
	Bestimmen Sie den Konvergenzradius und geben Sie den Konvergenzbereich der folgenden Reihen an. Untersuchen Sie dabei auch die R�nder des Konvergenzbereichs.
	\begin{eqnarray*}
		\begin{array}{l l l}
			(a)\; P(x) = \displaystyle{\sum\limits_{k=1}^\infty} k\cdot x^k \quad & (b)\; P(x) = \displaystyle{\sum\limits_{k=1}^\infty} (-1)^k \cdot \displaystyle{\frac{x^k}{k}} \quad & (c)\; P(x) = \displaystyle{\sum\limits_{k=1}^\infty} \displaystyle{\frac{x^k}{k^2}} \\
			(d)\; P(x) = \displaystyle{\sum\limits_{k=0}^\infty} \displaystyle{\frac{x^k}{2^k}} \quad & (e)\; P(x) = \displaystyle{\sum\limits_{k=0}^\infty } \displaystyle{\frac{k}{k+1} x^{k+1}} \quad & (f)\; P(x) = \displaystyle{\sum\limits_{k=0}^\infty} \displaystyle{\frac{k+1}{k!}} x^k
		\end{array}
	\end{eqnarray*}
\hspace{0.3cm}
\newline
\textbf{L�sung:}
Diese Aufgabe entspricht Aufgabe \(1\) in Papula Band 1, Abschnitt VI.2. Die Aufgabenstellung finden Sie auf Seite \(635\) und die detaillierten L�sungen ab Seite \(810\).
	\begin{itemize}
		\item[(a)] $R=1$ und R�nder $x=\pm 1$ divergent: Konvergenzbereich $(-1,1)$
		\item[(b)] $R=1$ und R�nder $x=- 1$ divergent (harmonische Reihe) und $x=1$ konvergent (Leibnitz) Konvergenzbereich $(-1,1]$
		\item[(c)] $R=1$ und R�nder $x=\pm 1$ konvergnt: Konvergenzbereich $[-1,1]$
		\item[(d)] $R=2$ und R�nder $x=\pm 2$ divergent: Konvergenzbereich $(-2,2)$
		\item[(e)] $R=1$ und R�nder $x=\pm 1$ divergent, weil notwendiges Kriterium f�r Konvergenz \( (a_k)_k\) Nullfolge nicht erf�llt ist: Konvergenzbereich $(-1,1)$
		\item[(f)] $R=\infty$: Konvergenzbereich $\mathbb R$
	\end{itemize}
\end{exercisebox}

%Moodle %HMI I, Blatt 11, Aufgabe 1 Moodle
\begin{exercisebox}[Konvergenz] %HMI I, Blatt 10, Aufgabe 4 Moodle
	Bestimmen Sie den Konvergenzradius der folgenden Reihen.
	\begin{eqnarray*}
		(a)\; \displaystyle{\sum\limits_{k=1}^\infty} \left(  \displaystyle{\frac{k!}{3 \cdot 5\cdot \dots \cdot (2k+1)}}\right)^{2} x^k \quad  (b)\; \displaystyle{\sum\limits_{k=1}^\infty} \left( \begin{array}{c} 2k\\ k\end{array}\right) x^k \quad  (c)\; \displaystyle{\sum\limits_{k=1}^\infty } \left( k^4 - 4k^3\right) x^k
	\end{eqnarray*}

\hspace{0.3cm}
\newline
\textbf{L�sung:}
	Konvergenz einer Potenzreihe liegt oft nicht f�r alle \(x\in \mathbb R\) vor\footnote{richtig klar wird das erst, wenn man Funktionen im Komplexen diskutiert. Dies ist Teil der mathematischen Disziplin Funktionentheorie, der wir uns leider nicht widmen werden}. Es stellt sich also die Frage: f�r welche \(x\) konvergiert eine Potenzreihe. Innerhalb des Konvergenzgebietes stellt die Potenzreihe eine Funktion dar!!! Das ist doch erstaunlich, oder? Wie auch immer: es gibt zwei wichtige Kriterien, um das Konverengzgebiet, den sogenannten Konvergenzradius \(R\) einer Potenzreihe zu bestimmen: das Quotientenkriterium und das Wurzelkriterium. Es gilt
	\begin{eqnarray*}
		\textnormal{Quotientenkriterium}: \quad &&\lim\limits_{n\to \infty} \left| \frac{ a_{n+1}}{a_n}\right| = a < \infty \Rightarrow R=\frac1a\,,\\
		\textnormal{Wurzelkriterium}: \quad &&\lim\limits_{n\to \infty} \sqrt[n]{\left| a_n\right|} = a < \infty \Rightarrow R=\frac1a\,.
	\end{eqnarray*}
	\begin{itemize}
		\item[(a)] Wir bestimmen den Konvergenzradius mit Hilfe des Quotientenkriteriums.
			\begin{eqnarray*} &&\lim\limits_{n\to \infty} \left| \frac{ ((n+1)!)^2}{\left(3\cdot 5 \cdot \cdots\cdot (2n+1)\cdot (2(n+1)+1) \right)^2} \cdot \frac{\left( 3\cdot 5 \cdot \cdots \cdot (2n+1)\right)^2}{(n!)^2 } \right| \\
				&=& \lim\limits_{n\to \infty} \left| \frac{n^2+2n +1 }{4 n^2 + 12 n + 9}\right| = \lim\limits_{n\to \infty} \left| \frac{1+2/n +1/n^2 }{4 + 12/ n + 9/n^2}\right|  \\ &=& \frac14 \Rightarrow R = 4\end{eqnarray*}
		\item[(b)] Wir bestimmen den Konvergenzradius mit Hilfe des Wurzelkriteriums.
			\begin{eqnarray*} \lim\limits_{n\to \infty} \sqrt[n]{{2n}\choose{n}} =  4 \Rightarrow R = \frac14\end{eqnarray*}
		\item[(c)] Wir bestimmen den Konvergenzradius mit Hilfe des Quotientenkriteriums.
			\begin{eqnarray*} \lim\limits_{n\to \infty} \left| \frac{ (n+1)^3(n+1-4)}{(n^4-4n^3)} \right| = \lim\limits_{n\to \infty} \left| \frac{n^3\cdot (1+2/n - \cdots) }{n^3\cdot (1-4/n)}\right| = 1 \Rightarrow R = 1\end{eqnarray*}
	\end{itemize}
\end{exercisebox}

\begin{exercisebox}[ Analytische Funktionen]
Begr�nden Sie, warum die folgende Funktion nicht analytisch ist
\[ f: x\mapsto  \begin{cases} e^{-1/x^2}, \quad x\not=0\,,\\
0, \quad x=0\,.\end{cases}
\]

\hspace{0.3cm}
\newline
\textbf{L�sung:}
Diese  Funktion  schmiegt  sich  bei $x=0$ extrem  stark  an  die $x$-Achse:  Der Funktionswert und alle Ableitungen sind dort $0$. Die Taylor-Reihe dieser Funktion $f$ f�r $x_0=0$ ist deshalb f�r alle $x$ gleich null - also nicht gleich $f$.
\end{exercisebox}

\begin{exercisebox}[Leibnitz-Kriterium] %HMI I, Blatt 10, Aufgabe 2
	Der Weg eines Teilchens in Abh�ngigkeit der Beschleunigung $x<0$ sei durch die unendliche Reihe
	\[
		S(x) = \sum\limits_{k=1}^\infty \frac{x^{k-1}}{k^2\cdot 3^k}
	\] gegeben.  Untersuchen Sie mit Hilfe des Leibnizkriteriums f�r welche $x\in\mathbb R$ die Reihe $S(x)$ sinnvoll gegeben ist.
%Geben Sie eine hinreichende Bedingung f�r die Beschleunigung an, so dass $0\leq S(x) \leq \frac13$ gilt.

\hspace{0.3cm}
\newline
\textbf{L�sung:}
\(S(x)\) ist sinnvoll definiert, falls die Reihe konvergiert:
			\[ \sum\limits_{k=0}^\infty \frac{x^{k-1}}{k^2\cdot 3^k} = \sum\limits_{k=1}^\infty (-1)^{k-1}\frac{|x|^{k-1}}{n^2\cdot 3^k}\]
			Das hei�t, dass wir das Leibnitz-Kriterium anwenden k�nnen. Eine Reihe �ber eine alternierende, monoton fallende Nullfolge ist nach Leibnitz konvergent. Die Frage ist also: f�r welche \(x\) ist die Folge \( a_k:= \displaystyle{ \frac{|x|^{k-1}}{k^2\cdot 3^k} } \)  monton fallend?
			\begin{itemize}
				\item offensichtlich ist Null eine untere Schranke f�r \(a_k\): \(0 < a_k, \; k\in \mathbb N^+\)
				\item die Folge f�llt monoton, wenn der Quotient aus zwei Folgegliedern kleiner als eins ist, also
					\begin{eqnarray*} \left| \frac{a_{k+1}}{a_k} \right| &=&  \frac{ |x|^k \cdot k^2 \cdot 3^k}{(k+1)^2 \cdot 3^{k+1} \cdot |x|^{k+1}}\\
						&=&  \frac{|x|}{3}\left( \frac{n}{n+1}\right)^2 \stackrel{!}{\leq} 1  \Leftrightarrow |x| \leq 3 \end{eqnarray*}
				\end{itemize}
						Also ist \(S(x)\) wohldefiniert f�r \(-3 \leq x < 0\).
\end{exercisebox}

\begin{exercisebox}[Potenzreihen und Differentialgleichungen] %Alternativ: Blatt 11, Aufgabe 3
  Ein K�rper mit der Masse $m$ bewege sich auf einer geradlinigen Bahn. Die Anfangsgeschwindigkeit sei $v_0 = 1$. Durch eine Kraft $F$, die der Geschwindigkeit proportional ist, wird der K�rper abgebremst.

  \begin{itemize}
	  \item[(a)] Stellen Sie  Bewegungsgleichung f�r den K�rper auf.
	  \item[(b)] L�sen Sie die Differentialgleichung mit Hilfe eines Potenzreihenansatzes.
  \end{itemize}

\hspace{0.3cm}
\newline
\textbf{L�sung:}
  \begin{itemize}
	  \item[(a)]
  In dieser Aufgabe ist eine Bewegungsgleichung zu entwickeln, wie Sie sie aus der Physik kennen. Zu jeder Zeit gilt $F_c(t) = F_a(t)$ mit
  \begin{eqnarray*}
	  F_c(t) &=& -c \cdot v(t) \quad \textnormal{(Bremskraft)} \\
  F_m(t) &=& m \cdot a(t) \quad  \textnormal{(Newton)} \,.
  \end{eqnarray*}
  Mit \[  a(t) = \dot v(t)\,,\] folgt ein sogenanntes Anfangswertproblem f�r $v$, das ist eine Differentialgleichung f�r $v$ mit gegebener Angfangsgeschwindigkeit $ v_0$:
  \[ \begin{cases}
		  \dot v(t) &= -\frac{c}{m} v(t) \quad t > 0\,, \quad (*)\\
		  v(0) = v_0 &= 1\,.
	  \end{cases}
  \]
  \item[(b)]
  Die Idee ist diese Differentialgleichung mit einem Potenzreihenansatz zu l�sen. Wir nehmen an \[v(t) = \sum\limits_{n=0}^\infty a_n t^n \Rightarrow \dot v(t) = \sum\limits_{n=0}^{\infty} (n+1)\cdot a_{n+1} t^{n}\,.\]
  Setzen wir nun die Reihenentwicklungen f�r $v$ und $\dot v$ in $(*)$ ein, dann
  \begin{eqnarray*}
	  \sum\limits_{n=0}^{\infty} (n+1)\cdot a_{n+1} t^{n} &=& -\frac{c}{m} \cdot \sum\limits_{n=0}^\infty a_n t^n \\
	  \Rightarrow && \sum\limits_{n=0}^\infty \left( (n+1)\cdot a_{n+1} + \frac{c}{m} a_{n} \right) t^n = 0   \quad t> 0 \\
		   \sum\limits_{n=0}^\infty a_n 0^n = 1 \Rightarrow a_0 = 1\,.
	   \end{eqnarray*}
	   Das hei�t, dass die  unbekannten Koeffizienten  $a_n, n\geq2$ rekursiv miteinander verkn�pft sind:
  \begin{eqnarray*}
	  a_0 &=& 1 \\
	  a_{n+1} = -\frac{c}{m} \frac{1}{n+1} a_n \Rightarrow a_1 = -\frac{c}{m},\; a_2 =  \left( -\frac{c}{m}\right)^2 \frac{1}{2},\; a_3 =  \left( -\frac{c}{m}\right)^3 \frac{1}{2\cdot 3} \cdots
	   \end{eqnarray*}
	   Mit vollst�ndiger Induktion beweist man \[ a_n = \left( -\frac{m}{c}\right)^n \frac{1}{n!} \,.\]
	   Das hei�t wir haben die L�sung in Form einer Potenzreihe gefunden (deren Wert kennen wir auch schon!):
	   \[ v(t) = \sum\limits_{n=0}^\infty  \left(  -\frac{c}{m} t\right)^n \frac{1}{n!} \Rightarrow v(t) = e^{-m/c t} \,.\]
\end{itemize}
\end{exercisebox}

